\chapter{Reproducing the Results}
\label{app:repro}

Every metric in Chapter~\ref{ch:results} is regenerated from committed fingerprints by a command in \file{control/cmd}, and a test compares each command's output with the committed metric file, so a change that alters a reported number fails CI. Run from \file{control/}:

\begin{Verbatim}[fontsize=\footnotesize,frame=single,framesep=2mm]
# divergence per cell pair and class (Table 11.3)
go run ./cmd/caliper-divergence-metrics \
    -out ../fixtures/metrics/divergence-metrics.json
# expectation, claim and undeclared findings at the restricted tier (Table 11.4)
go run ./cmd/caliper-ci-report -cell-fixtures ../fixtures/cells \
    -metrics-out ../fixtures/metrics/finding-metrics.json
# attribution against single-cause environments (Table 11.7)
go run ./cmd/caliper-attribution-accuracy \
    -ground-truth ../fixtures/ground-truth/attribution.json \
    -out ../fixtures/metrics/attribution-accuracy.json
# closure verified by re-probing (Table 11.9)
go run ./cmd/caliper-closure-metrics \
    -out ../fixtures/metrics/closure-metrics.json
# residual state and runtime (Table 11.10)
go run ./cmd/caliper-residual-metrics \
    -out ../fixtures/metrics/residual-metrics.json
\end{Verbatim}

The measurements themselves come from the cells, with \file{deploy/gce/cells.sh}:

\begin{Verbatim}[fontsize=\footnotesize,frame=single,framesep=2mm]
# a cell: create (VM from a pinned image, then provisioning), measure, stop
deploy/gce/cells.sh create  cellN
deploy/gce/cells.sh measure cellN <image> > cellN.raw   # reboots first
deploy/gce/cells.sh stop    cellN                       # a running VM is billed

# ground-truth modes and the hand-hardened profile on cell 2
deploy/gce/cells.sh measure cell2 <image> "" <mode>
deploy/gce/cells.sh measure cell2 <image> \
    engine/fixtures/seccomp/hardened.json

# reduced corpus; closure pipeline; sandbox handlers with seccomp
CALIPER_CORPUS_SET=reduced deploy/gce/cells.sh measure cellN <image>
CALIPER_POSTURE_TIER=untrusted-multi-tenant \
    deploy/gce/run-pipeline.sh cellN
CALIPER_HANDLER=runsc-seccomp deploy/gce/run-pipeline.sh cell5
CALIPER_HANDLER=kata-seccomp  deploy/gce/run-pipeline.sh cell6

# complete a raw measurement into a fingerprint (from control/)
go run ./cmd/caliper-complete -distribution "Ubuntu 24.04" \
    -runtime containerd -runtime-version 2.2.1 -runtime-class runc \
    -node-kernel 7.0.0-1014-gcp < ../cell2.raw > ../fixtures/cells/cell2.json
\end{Verbatim}

The released image is verified with cosign against the release workflow's identity (\file{RELEASE.md}):

\begin{Verbatim}[fontsize=\footnotesize,frame=single,framesep=2mm]
ghcr.io/abhinavmadake/caliper@sha256:
    61b918d61aae4cec0b9df70526555c65ff42f9a33f6381ebb99eb1f83a174db7
\end{Verbatim}

The production undecidable fractions of Table~\ref{tab:undecidable} are not a committed metric. They were computed for this report by calling \texttt{control.Attribute} and \texttt{control.BuildReport} on the baseline fingerprints in \file{fixtures/pipeline/} and on \file{fixtures/gke-cos-containerd-2.2.7.json}; for a single fingerprint this is the value the text report prints on its ``Undecidable'' line.

\chapter{Project Record}
\label{app:record}

\section{Timeline}

The repository's first commit is dated 2026-09-04 and week 1 began on Monday 2026-09-07. Figure~\ref{fig:gantt} shows the main periods of work, taken from the merge dates in the repository's history, and Figure~\ref{fig:commits} the commits on \texttt{main} per calendar week (each merged pull request is one squashed commit).

\begin{figure}[H]
\centering
\begin{tikzpicture}[x=0.165cm,y=0.52cm]
% day 0 = 2026-09-04; ticks every 7 days from 2026-09-07 (day 3)
\foreach \d/\l in {3/{Sep 7},10/{Sep 14},17/{Sep 21},24/{Sep 28},31/{Oct 5},38/{Oct 12},45/{Oct 19},52/{Oct 26},59/{Nov 2},66/{Nov 9}} {
  \draw[gridc] (\d,0.4) -- (\d,-11.6);
  \node[font=\scriptsize,text=inkm,anchor=south] at (\d,0.4) {\l};
}
\newcommand{\gbar}[4]{\fill[#4] (#1,-#3+0.32) rectangle (#2,-#3-0.32);}
\newcommand{\glab}[2]{\node[font=\scriptsize,anchor=east,text=inkm] at (-0.5,-#2) {#1};}
\glab{Interfaces (spec)}{1}            \gbar{0}{7}{1}{s1}
\glab{Engine, reference corpus}{2}     \gbar{6}{10}{2}{s1}
\glab{Corpus families}{3}              \gbar{10}{23}{3}{s1}
\glab{Fingerprint, diff, classifier}{4}\gbar{10}{45}{4}{s1}
\glab{Evaluator, posture, reporting}{5}\gbar{36}{42}{5}{s1}
\glab{Attribution and accuracy}{6}     \gbar{35}{52}{6}{s1}
\glab{GKE CI, six cells}{7}            \gbar{42}{44}{7}{s1}
\glab{Remediation, SPO}{8}             \gbar{44}{50}{8}{s1}
\glab{Evaluation metrics}{9}           \gbar{48}{53}{9}{s1}
\glab{Release v1.0.0}{10}              \node[circle,fill=s2,inner sep=0pt,minimum size=5pt] at (55,-10) {};
\glab{Disclosure review, report}{11}   \gbar{63}{65}{11}{s1}
\end{tikzpicture}
\caption{Main periods of work (from merge dates). The week-17 gate, with the demonstration, is 2026-11-29.}
\label{fig:gantt}
\end{figure}

\begin{figure}[H]
\centering
\begin{tikzpicture}
\begin{axis}[cal, ybar, width=0.8\textwidth, height=5.0cm, bar width=14pt,
  ymin=0, ymax=36, ylabel={Commits on \texttt{main}},
  symbolic x coords={Sep 7,Sep 14,Sep 21,Sep 28,Oct 5,Oct 12,Oct 19,Oct 26},
  xtick=data, ymajorgrids, xlabel={Week beginning},
  nodes near coords, every node near coord/.append style={font=\scriptsize,text=inkm}]
\addplot[fill=s1,draw=none] coordinates {(Sep 7,30) (Sep 14,18) (Sep 21,4) (Sep 28,3) (Oct 5,14) (Oct 12,23) (Oct 19,13) (Oct 26,11)};
\end{axis}
\end{tikzpicture}
\caption{Commits on \texttt{main} per week, 116 in all (the first week includes the three commits of 2026-09-04).}
\label{fig:commits}
\end{figure}

\section{Milestones}

The proposal's schedule was tracked as one GitHub milestone per freeze gate (Table~\ref{tab:milestones}). Week 16 was kept empty as declared slack; Table~\ref{tab:gates} in Chapter~\ref{ch:proposed} compares each gate's due date with the date it was met.


\begin{table}[H]
\centering\small
\begin{tabular}{@{}llr@{}}
\toprule
\textbf{Milestone} & \textbf{Due} & \textbf{Issues} \\
\midrule
W2: interfaces frozen & 2026-09-11 & 2 \\
W3: engine and reference corpus & 2026-09-18 & 10 \\
W8: corpus frozen & 2026-09-27 & 14 \\
W9: fingerprint frozen & 2026-10-11 & 14 \\
W12: attribution, evaluator, remediation & 2026-10-25 & 21 \\
W15: evaluation complete & 2026-11-15 & 9 \\
W16: integration buffer (kept empty) & 2026-11-22 & 0 \\
W17: release & 2026-11-29 & 3 \\
\bottomrule
\end{tabular}
\caption{Milestones and their issues. The backlog was seeded from the proposal by \file{scripts/seed-backlog.sh}.}
\label{tab:milestones}
\end{table}

\section{Statistics}

At the time of writing \texttt{main} has 116 commits from 87 merged pull requests, and 71 of 73 issues are closed; the two open ones are this report and the demonstration. The code is about 9,900 lines of Rust, 8,900 lines of Go and 2,400 lines of shell, tests included. The repository holds 8 fingerprints for the cells and the GKE node, 7 ground-truth fingerprints, 8 pipeline runs, 13 runtime runs and 5 metric files. Workstreams B and D exchanged their work on 2026-09-14, after the week-3 gate; the letters stayed with the work.

\chapter{The Complete Probe Index}
\label{app:index}

Table~\ref{tab:index} lists every probe in \file{corpus/index.json}: the 185 committed probes by family, then the 28 deferred probes. ``Isolates'' is the oracle's isolation (\texttt{seccomp}; \texttt{sec./LSM} for seccomp or an LSM; \texttt{undec.} for undecidable from the errno alone). ``G'' is the errno the kernel returns if the call reaches the oracle's point, and ``A'' the errno that absence of the entry point produces, with the encoded absence after a slash for io\_uring opcodes. Errnos are raw numbers, as in the fingerprint: 2 \errno{ENOENT}, 7 \errno{E2BIG}, 8 \errno{ENOEXEC}, 9 \errno{EBADF}, 14 \errno{EFAULT}, 19 \errno{ENODEV}, 22 \errno{EINVAL}, 38 \errno{ENOSYS}, 61 \errno{ENODATA}, 93 \errno{EPROTONOSUPPORT}, 95 \errno{EOPNOTSUPP}, 97 \errno{EAFNOSUPPORT}, 126 \errno{ENOKEY}. Every committed probe has a 1,000 ms timeout and runs in a forked child.

\begingroup\scriptsize\setlength{\tabcolsep}{3pt}\renewcommand{\arraystretch}{1.0}
\begin{longtable}{@{}P{4.4cm}P{1.35cm}R{0.75cm}R{0.95cm}P{7.65cm}@{}}
\caption{The complete probe index (\file{corpus/index.json}). G: guaranteed errno; A: absence errno (encoded absence after a slash); $\dagger$: may autoload a module.}\label{tab:index}\\
\toprule \textbf{Probe} & \textbf{Isolates} & \textbf{G} & \textbf{A} & \textbf{Operation} \\ \midrule\endfirsthead
\multicolumn{5}{@{}l}{\itshape Table \thetable{} continued}\\ \toprule \textbf{Probe} & \textbf{Isolates} & \textbf{G} & \textbf{A} & \textbf{Operation} \\ \midrule\endhead
\bottomrule\endfoot
\multicolumn{5}{@{}l}{\textbf{\texttt{socket}} (24)}\\
\pid{socket.family.af_inet} & sec./LSM & -- & -- & socket(AF\_\allowbreak{}INET, SOCK\_\allowbreak{}STREAM, 0): the family every workload uses \\
\pid{socket.family.af_alg}$^{\dagger}$ & sec./LSM & -- & 97 & socket(AF\_\allowbreak{}ALG, SOCK\_\allowbreak{}SEQPACKET, 0): the crypto user API is reachable \\
\pid{socket.family.af_packet} & undec. & -- & 97 & socket(AF\_\allowbreak{}PACKET, SOCK\_\allowbreak{}RAW, 0): raw link-layer access, CAP\_\allowbreak{}NET\_\allowbreak{}RAW \\
\pid{socket.family.out_of_range} & seccomp & 97 & -- & socket(1000, SOCK\_\allowbreak{}STREAM, 0): a family beyond NPROTO, guaranteed EAFNOSUPPORT \\
\pid{socket.family.af_unix} & sec./LSM & -- & -- & socket(AF\_\allowbreak{}UNIX, SOCK\_\allowbreak{}STREAM, 0): the local family every runtime needs \\
\pid{socket.family.af_inet6}$^{\dagger}$ & sec./LSM & -- & 97 & socket(AF\_\allowbreak{}INET6, SOCK\_\allowbreak{}STREAM, 0): IPv6 is reachable \\
\pid{socket.family.af_key}$^{\dagger}$ & undec. & -- & 97 & socket(AF\_\allowbreak{}KEY, SOCK\_\allowbreak{}RAW, PF\_\allowbreak{}KEY\_\allowbreak{}V2): the IPsec key management interface, CAP\_\allowbreak{}NET\_\allowbreak{}ADMIN \\
\pid{socket.family.af_bluetooth}$^{\dagger}$ & sec./LSM & -- & 97 & socket(AF\_\allowbreak{}BLUETOOTH, SOCK\_\allowbreak{}STREAM, BTPROTO\_\allowbreak{}RFCOMM): the Bluetooth stack is reachable. Only from the host network namespace: in any other, bt\_\allowbreak{}sock\_\allowbreak{}create answers EAFNOSUPPORT whatever the policy \\
\pid{socket.family.af_can}$^{\dagger}$ & sec./LSM & -- & 97 & socket(AF\_\allowbreak{}CAN, SOCK\_\allowbreak{}RAW, CAN\_\allowbreak{}RAW): the CAN bus stack is reachable \\
\pid{socket.family.af_tipc}$^{\dagger}$ & sec./LSM & -- & 97 & socket(AF\_\allowbreak{}TIPC, SOCK\_\allowbreak{}RDM, 0): the TIPC cluster stack is reachable \\
\pid{socket.family.af_rds}$^{\dagger}$ & sec./LSM & -- & 97 & socket(AF\_\allowbreak{}RDS, SOCK\_\allowbreak{}SEQPACKET, 0): Reliable Datagram Sockets are reachable \\
\pid{socket.family.af_xdp} & undec. & -- & 97 & socket(AF\_\allowbreak{}XDP, SOCK\_\allowbreak{}RAW, 0): AF\_\allowbreak{}XDP zero-copy sockets, CAP\_\allowbreak{}NET\_\allowbreak{}RAW \\
\pid{socket.family.af_rxrpc}$^{\dagger}$ & sec./LSM & -- & 97 & socket(AF\_\allowbreak{}RXRPC, SOCK\_\allowbreak{}DGRAM, PF\_\allowbreak{}INET): the RxRPC transport is reachable \\
\pid{socket.family.af_phonet}$^{\dagger}$ & undec. & -- & 97 & socket(AF\_\allowbreak{}PHONET, SOCK\_\allowbreak{}DGRAM, 0): the Phonet stack, CAP\_\allowbreak{}SYS\_\allowbreak{}ADMIN \\
\pid{socket.family.af_ieee802154}$^{\dagger}$ & sec./LSM & -- & 97 & socket(AF\_\allowbreak{}IEEE802154, SOCK\_\allowbreak{}DGRAM, 0): the 802.15.4 stack is reachable \\
\pid{socket.family.af_nfc}$^{\dagger}$ & sec./LSM & -- & 97 & socket(AF\_\allowbreak{}NFC, SOCK\_\allowbreak{}STREAM, NFC\_\allowbreak{}SOCKPROTO\_\allowbreak{}LLCP): the NFC stack is reachable. Only from the host network namespace: in any other, nfc\_\allowbreak{}sock\_\allowbreak{}create answers EAFNOSUPPORT whatever the policy \\
\pid{socket.family.af_vsock}$^{\dagger}$ & sec./LSM & -- & 97 & socket(AF\_\allowbreak{}VSOCK, SOCK\_\allowbreak{}STREAM, 0): the host-guest vsock family is reachable \\
\pid{socket.family.af_kcm}$^{\dagger}$ & sec./LSM & -- & 97 & socket(AF\_\allowbreak{}KCM, SOCK\_\allowbreak{}DGRAM, KCMPROTO\_\allowbreak{}CONNECTED): kernel connection multiplexor \\
\pid{socket.family.af_qipcrtr}$^{\dagger}$ & sec./LSM & -- & 97 & socket(AF\_\allowbreak{}QIPCRTR, SOCK\_\allowbreak{}DGRAM, PF\_\allowbreak{}QIPCRTR): the Qualcomm IPC router is reachable \\
\pid{socket.family.af_smc}$^{\dagger}$ & sec./LSM & -- & 97 & socket(AF\_\allowbreak{}SMC, SOCK\_\allowbreak{}STREAM, SMCPROTO\_\allowbreak{}SMC): shared memory communications \\
\pid{socket.family.af_mctp}$^{\dagger}$ & sec./LSM & -- & 97 & socket(AF\_\allowbreak{}MCTP, SOCK\_\allowbreak{}DGRAM, 0): the Management Component Transport Protocol \\
\pid{socket.family.af_ax25}$^{\dagger}$ & sec./LSM & -- & 97 & socket(AF\_\allowbreak{}AX25, SOCK\_\allowbreak{}DGRAM, 0): the amateur-radio AX.25 stack is reachable \\
\pid{socket.family.af_appletalk}$^{\dagger}$ & sec./LSM & -- & 97 & socket(AF\_\allowbreak{}APPLETALK, SOCK\_\allowbreak{}DGRAM, 0): the AppleTalk stack is reachable \\
\pid{socket.family.af_x25}$^{\dagger}$ & sec./LSM & -- & 97 & socket(AF\_\allowbreak{}X25, SOCK\_\allowbreak{}SEQPACKET, 0): the X.25 stack is reachable \\
\multicolumn{5}{@{}l}{\textbf{\texttt{netlink}} (18)}\\
\pid{netlink.protocol.route} & sec./LSM & -- & -- & socket(AF\_\allowbreak{}NETLINK, SOCK\_\allowbreak{}RAW, NETLINK\_\allowbreak{}ROUTE): rtnetlink is reachable \\
\pid{netlink.protocol.usersock} & sec./LSM & -- & -- & socket(AF\_\allowbreak{}NETLINK, SOCK\_\allowbreak{}RAW, NETLINK\_\allowbreak{}USERSOCK): the user-space-to-user-space protocol \\
\pid{netlink.protocol.sock_diag} & sec./LSM & -- & 93 & socket(AF\_\allowbreak{}NETLINK, SOCK\_\allowbreak{}RAW, NETLINK\_\allowbreak{}SOCK\_\allowbreak{}DIAG): socket enumeration (ss, inet\_\allowbreak{}diag) \\
\pid{netlink.protocol.nflog}$^{\dagger}$ & sec./LSM & -- & 93 & socket(AF\_\allowbreak{}NETLINK, SOCK\_\allowbreak{}RAW, NETLINK\_\allowbreak{}NFLOG): the legacy netfilter log protocol \\
\pid{netlink.protocol.xfrm}$^{\dagger}$ & sec./LSM & -- & 93 & socket(AF\_\allowbreak{}NETLINK, SOCK\_\allowbreak{}RAW, NETLINK\_\allowbreak{}XFRM): IPsec transform configuration \\
\pid{netlink.protocol.selinux} & sec./LSM & -- & 93 & socket(AF\_\allowbreak{}NETLINK, SOCK\_\allowbreak{}RAW, NETLINK\_\allowbreak{}SELINUX): SELinux event notifications \\
\pid{netlink.protocol.iscsi}$^{\dagger}$ & sec./LSM & -- & 93 & socket(AF\_\allowbreak{}NETLINK, SOCK\_\allowbreak{}RAW, NETLINK\_\allowbreak{}ISCSI): the open-iSCSI transport \\
\pid{netlink.protocol.audit} & sec./LSM & -- & 93 & socket(AF\_\allowbreak{}NETLINK, SOCK\_\allowbreak{}RAW, NETLINK\_\allowbreak{}AUDIT): the audit subsystem \\
\pid{netlink.protocol.fib_lookup} & sec./LSM & -- & 93 & socket(AF\_\allowbreak{}NETLINK, SOCK\_\allowbreak{}RAW, NETLINK\_\allowbreak{}FIB\_\allowbreak{}LOOKUP): route lookups from user space \\
\pid{netlink.protocol.connector}$^{\dagger}$ & sec./LSM & -- & 93 & socket(AF\_\allowbreak{}NETLINK, SOCK\_\allowbreak{}RAW, NETLINK\_\allowbreak{}CONNECTOR): the kernel connector (proc events) \\
\pid{netlink.protocol.netfilter}$^{\dagger}$ & sec./LSM & -- & 93 & socket(AF\_\allowbreak{}NETLINK, SOCK\_\allowbreak{}RAW, NETLINK\_\allowbreak{}NETFILTER): nfnetlink, the nf\_\allowbreak{}tables entry point \\
\pid{netlink.protocol.kobject_uevent} & sec./LSM & -- & 93 & socket(AF\_\allowbreak{}NETLINK, SOCK\_\allowbreak{}RAW, NETLINK\_\allowbreak{}KOBJECT\_\allowbreak{}UEVENT): device hotplug events \\
\pid{netlink.protocol.generic} & sec./LSM & -- & -- & socket(AF\_\allowbreak{}NETLINK, SOCK\_\allowbreak{}RAW, NETLINK\_\allowbreak{}GENERIC): generic netlink, every modern family's bus \\
\pid{netlink.protocol.scsitransport} & sec./LSM & -- & 93 & socket(AF\_\allowbreak{}NETLINK, SOCK\_\allowbreak{}RAW, NETLINK\_\allowbreak{}SCSITRANSPORT): the SCSI transport netlink \\
\pid{netlink.protocol.rdma}$^{\dagger}$ & sec./LSM & -- & 93 & socket(AF\_\allowbreak{}NETLINK, SOCK\_\allowbreak{}RAW, NETLINK\_\allowbreak{}RDMA): the RDMA subsystem's netlink \\
\pid{netlink.protocol.crypto}$^{\dagger}$ & sec./LSM & -- & 93 & socket(AF\_\allowbreak{}NETLINK, SOCK\_\allowbreak{}RAW, NETLINK\_\allowbreak{}CRYPTO): the crypto user configuration API \\
\pid{netlink.protocol.smc}$^{\dagger}$ & sec./LSM & -- & 93 & socket(AF\_\allowbreak{}NETLINK, SOCK\_\allowbreak{}RAW, NETLINK\_\allowbreak{}SMC): the SMC protocol's netlink \\
\pid{netlink.protocol.out_of_range} & sec./LSM & 93 & -- & socket(AF\_\allowbreak{}NETLINK, SOCK\_\allowbreak{}RAW, 32): a protocol at MAX\_\allowbreak{}LINKS, guaranteed EPROTONOSUPPORT \\
\multicolumn{5}{@{}l}{\textbf{\texttt{io\_uring}} (53)}\\
\pid{io_uring.register.probe} & seccomp & 9 & 38 & io\_\allowbreak{}uring\_\allowbreak{}register(INT\_\allowbreak{}MAX, IORING\_\allowbreak{}REGISTER\_\allowbreak{}PROBE, NULL, 0): guaranteed EBADF \\
\pid{io_uring.setup.zero_entries} & seccomp & 22 & 38 & io\_\allowbreak{}uring\_\allowbreak{}setup(0, \&params): guaranteed EINVAL, no ring is created \\
\pid{io_uring.opcode.nop} & seccomp & -- & 38/61 & IORING\_\allowbreak{}REGISTER\_\allowbreak{}PROBE reports IORING\_\allowbreak{}OP\_\allowbreak{}NOP (0, since 5.1): the control: a ring that can be probed at all supports it \\
\pid{io_uring.opcode.readv} & seccomp & -- & 38/61 & IORING\_\allowbreak{}REGISTER\_\allowbreak{}PROBE reports IORING\_\allowbreak{}OP\_\allowbreak{}READV (1, since 5.1): vectored reads without read(2) \\
\pid{io_uring.opcode.writev} & seccomp & -- & 38/61 & IORING\_\allowbreak{}REGISTER\_\allowbreak{}PROBE reports IORING\_\allowbreak{}OP\_\allowbreak{}WRITEV (2, since 5.1): vectored writes without write(2) \\
\pid{io_uring.opcode.fsync} & seccomp & -- & 38/61 & IORING\_\allowbreak{}REGISTER\_\allowbreak{}PROBE reports IORING\_\allowbreak{}OP\_\allowbreak{}FSYNC (3, since 5.1): fsync without fsync(2) \\
\pid{io_uring.opcode.read_fixed} & seccomp & -- & 38/61 & IORING\_\allowbreak{}REGISTER\_\allowbreak{}PROBE reports IORING\_\allowbreak{}OP\_\allowbreak{}READ\_\allowbreak{}FIXED (4, since 5.1): reads into registered buffers \\
\pid{io_uring.opcode.write_fixed} & seccomp & -- & 38/61 & IORING\_\allowbreak{}REGISTER\_\allowbreak{}PROBE reports IORING\_\allowbreak{}OP\_\allowbreak{}WRITE\_\allowbreak{}FIXED (5, since 5.1): writes from registered buffers \\
\pid{io_uring.opcode.poll_add} & seccomp & -- & 38/61 & IORING\_\allowbreak{}REGISTER\_\allowbreak{}PROBE reports IORING\_\allowbreak{}OP\_\allowbreak{}POLL\_\allowbreak{}ADD (6, since 5.1): poll without poll(2) \\
\pid{io_uring.opcode.sync_file_range} & seccomp & -- & 38/61 & IORING\_\allowbreak{}REGISTER\_\allowbreak{}PROBE reports IORING\_\allowbreak{}OP\_\allowbreak{}SYNC\_\allowbreak{}FILE\_\allowbreak{}RANGE (8, since 5.2): sync\_\allowbreak{}file\_\allowbreak{}range without the syscall \\
\pid{io_uring.opcode.sendmsg} & seccomp & -- & 38/61 & IORING\_\allowbreak{}REGISTER\_\allowbreak{}PROBE reports IORING\_\allowbreak{}OP\_\allowbreak{}SENDMSG (9, since 5.3): sendmsg without sendmsg(2) \\
\pid{io_uring.opcode.recvmsg} & seccomp & -- & 38/61 & IORING\_\allowbreak{}REGISTER\_\allowbreak{}PROBE reports IORING\_\allowbreak{}OP\_\allowbreak{}RECVMSG (10, since 5.3): recvmsg without recvmsg(2) \\
\pid{io_uring.opcode.timeout} & seccomp & -- & 38/61 & IORING\_\allowbreak{}REGISTER\_\allowbreak{}PROBE reports IORING\_\allowbreak{}OP\_\allowbreak{}TIMEOUT (11, since 5.4): timers in the ring \\
\pid{io_uring.opcode.accept} & seccomp & -- & 38/61 & IORING\_\allowbreak{}REGISTER\_\allowbreak{}PROBE reports IORING\_\allowbreak{}OP\_\allowbreak{}ACCEPT (13, since 5.5): accept without accept4(2) \\
\pid{io_uring.opcode.connect} & seccomp & -- & 38/61 & IORING\_\allowbreak{}REGISTER\_\allowbreak{}PROBE reports IORING\_\allowbreak{}OP\_\allowbreak{}CONNECT (16, since 5.5): connect without connect(2) \\
\pid{io_uring.opcode.fallocate} & seccomp & -- & 38/61 & IORING\_\allowbreak{}REGISTER\_\allowbreak{}PROBE reports IORING\_\allowbreak{}OP\_\allowbreak{}FALLOCATE (17, since 5.6): fallocate without the syscall \\
\pid{io_uring.opcode.openat} & seccomp & -- & 38/61 & IORING\_\allowbreak{}REGISTER\_\allowbreak{}PROBE reports IORING\_\allowbreak{}OP\_\allowbreak{}OPENAT (18, since 5.6): opening files without openat(2) \\
\pid{io_uring.opcode.close} & seccomp & -- & 38/61 & IORING\_\allowbreak{}REGISTER\_\allowbreak{}PROBE reports IORING\_\allowbreak{}OP\_\allowbreak{}CLOSE (19, since 5.6): close without close(2) \\
\pid{io_uring.opcode.statx} & seccomp & -- & 38/61 & IORING\_\allowbreak{}REGISTER\_\allowbreak{}PROBE reports IORING\_\allowbreak{}OP\_\allowbreak{}STATX (21, since 5.6): statx without the syscall \\
\pid{io_uring.opcode.read} & seccomp & -- & 38/61 & IORING\_\allowbreak{}REGISTER\_\allowbreak{}PROBE reports IORING\_\allowbreak{}OP\_\allowbreak{}READ (22, since 5.6): read without read(2) \\
\pid{io_uring.opcode.write} & seccomp & -- & 38/61 & IORING\_\allowbreak{}REGISTER\_\allowbreak{}PROBE reports IORING\_\allowbreak{}OP\_\allowbreak{}WRITE (23, since 5.6): write without write(2) \\
\pid{io_uring.opcode.fadvise} & seccomp & -- & 38/61 & IORING\_\allowbreak{}REGISTER\_\allowbreak{}PROBE reports IORING\_\allowbreak{}OP\_\allowbreak{}FADVISE (24, since 5.6): fadvise without the syscall \\
\pid{io_uring.opcode.madvise} & seccomp & -- & 38/61 & IORING\_\allowbreak{}REGISTER\_\allowbreak{}PROBE reports IORING\_\allowbreak{}OP\_\allowbreak{}MADVISE (25, since 5.6): madvise without madvise(2) \\
\pid{io_uring.opcode.send} & seccomp & -- & 38/61 & IORING\_\allowbreak{}REGISTER\_\allowbreak{}PROBE reports IORING\_\allowbreak{}OP\_\allowbreak{}SEND (26, since 5.6): send without send(2) \\
\pid{io_uring.opcode.recv} & seccomp & -- & 38/61 & IORING\_\allowbreak{}REGISTER\_\allowbreak{}PROBE reports IORING\_\allowbreak{}OP\_\allowbreak{}RECV (27, since 5.6): recv without recv(2) \\
\pid{io_uring.opcode.openat2} & seccomp & -- & 38/61 & IORING\_\allowbreak{}REGISTER\_\allowbreak{}PROBE reports IORING\_\allowbreak{}OP\_\allowbreak{}OPENAT2 (28, since 5.6): openat2 without the syscall \\
\pid{io_uring.opcode.epoll_ctl} & seccomp & -- & 38/61 & IORING\_\allowbreak{}REGISTER\_\allowbreak{}PROBE reports IORING\_\allowbreak{}OP\_\allowbreak{}EPOLL\_\allowbreak{}CTL (29, since 5.6): epoll\_\allowbreak{}ctl without the syscall \\
\pid{io_uring.opcode.splice} & seccomp & -- & 38/61 & IORING\_\allowbreak{}REGISTER\_\allowbreak{}PROBE reports IORING\_\allowbreak{}OP\_\allowbreak{}SPLICE (30, since 5.7): splice without splice(2) \\
\pid{io_uring.opcode.tee} & seccomp & -- & 38/61 & IORING\_\allowbreak{}REGISTER\_\allowbreak{}PROBE reports IORING\_\allowbreak{}OP\_\allowbreak{}TEE (33, since 5.8): tee without tee(2) \\
\pid{io_uring.opcode.shutdown} & seccomp & -- & 38/61 & IORING\_\allowbreak{}REGISTER\_\allowbreak{}PROBE reports IORING\_\allowbreak{}OP\_\allowbreak{}SHUTDOWN (34, since 5.11): shutdown without shutdown(2) \\
\pid{io_uring.opcode.renameat} & seccomp & -- & 38/61 & IORING\_\allowbreak{}REGISTER\_\allowbreak{}PROBE reports IORING\_\allowbreak{}OP\_\allowbreak{}RENAMEAT (35, since 5.11): renames without renameat2(2) \\
\pid{io_uring.opcode.unlinkat} & seccomp & -- & 38/61 & IORING\_\allowbreak{}REGISTER\_\allowbreak{}PROBE reports IORING\_\allowbreak{}OP\_\allowbreak{}UNLINKAT (36, since 5.11): unlinks without unlinkat(2) \\
\pid{io_uring.opcode.mkdirat} & seccomp & -- & 38/61 & IORING\_\allowbreak{}REGISTER\_\allowbreak{}PROBE reports IORING\_\allowbreak{}OP\_\allowbreak{}MKDIRAT (37, since 5.15): mkdir without mkdirat(2) \\
\pid{io_uring.opcode.symlinkat} & seccomp & -- & 38/61 & IORING\_\allowbreak{}REGISTER\_\allowbreak{}PROBE reports IORING\_\allowbreak{}OP\_\allowbreak{}SYMLINKAT (38, since 5.15): symlinks without symlinkat(2) \\
\pid{io_uring.opcode.linkat} & seccomp & -- & 38/61 & IORING\_\allowbreak{}REGISTER\_\allowbreak{}PROBE reports IORING\_\allowbreak{}OP\_\allowbreak{}LINKAT (39, since 5.15): hard links without linkat(2) \\
\pid{io_uring.opcode.msg_ring} & seccomp & -- & 38/61 & IORING\_\allowbreak{}REGISTER\_\allowbreak{}PROBE reports IORING\_\allowbreak{}OP\_\allowbreak{}MSG\_\allowbreak{}RING (40, since 5.18): messages and descriptors passed between rings \\
\pid{io_uring.opcode.fsetxattr} & seccomp & -- & 38/61 & IORING\_\allowbreak{}REGISTER\_\allowbreak{}PROBE reports IORING\_\allowbreak{}OP\_\allowbreak{}FSETXATTR (41, since 5.19): fsetxattr without the syscall \\
\pid{io_uring.opcode.setxattr} & seccomp & -- & 38/61 & IORING\_\allowbreak{}REGISTER\_\allowbreak{}PROBE reports IORING\_\allowbreak{}OP\_\allowbreak{}SETXATTR (42, since 5.19): setxattr without the syscall \\
\pid{io_uring.opcode.fgetxattr} & seccomp & -- & 38/61 & IORING\_\allowbreak{}REGISTER\_\allowbreak{}PROBE reports IORING\_\allowbreak{}OP\_\allowbreak{}FGETXATTR (43, since 5.19): fgetxattr without the syscall \\
\pid{io_uring.opcode.getxattr} & seccomp & -- & 38/61 & IORING\_\allowbreak{}REGISTER\_\allowbreak{}PROBE reports IORING\_\allowbreak{}OP\_\allowbreak{}GETXATTR (44, since 5.19): getxattr without the syscall \\
\pid{io_uring.opcode.socket} & seccomp & -- & 38/61 & IORING\_\allowbreak{}REGISTER\_\allowbreak{}PROBE reports IORING\_\allowbreak{}OP\_\allowbreak{}SOCKET (45, since 5.19): socket creation without socket(2), any family \\
\pid{io_uring.opcode.uring_cmd} & seccomp & -- & 38/61 & IORING\_\allowbreak{}REGISTER\_\allowbreak{}PROBE reports IORING\_\allowbreak{}OP\_\allowbreak{}URING\_\allowbreak{}CMD (46, since 5.19): driver passthrough commands (NVMe, ublk) without ioctl(2) \\
\pid{io_uring.opcode.send_zc} & seccomp & -- & 38/61 & IORING\_\allowbreak{}REGISTER\_\allowbreak{}PROBE reports IORING\_\allowbreak{}OP\_\allowbreak{}SEND\_\allowbreak{}ZC (47, since 6.0): zero-copy send \\
\pid{io_uring.opcode.sendmsg_zc} & seccomp & -- & 38/61 & IORING\_\allowbreak{}REGISTER\_\allowbreak{}PROBE reports IORING\_\allowbreak{}OP\_\allowbreak{}SENDMSG\_\allowbreak{}ZC (48, since 6.1): zero-copy sendmsg \\
\pid{io_uring.opcode.read_multishot} & seccomp & -- & 38/61 & IORING\_\allowbreak{}REGISTER\_\allowbreak{}PROBE reports IORING\_\allowbreak{}OP\_\allowbreak{}READ\_\allowbreak{}MULTISHOT (49, since 6.7): multishot reads \\
\pid{io_uring.opcode.waitid} & seccomp & -- & 38/61 & IORING\_\allowbreak{}REGISTER\_\allowbreak{}PROBE reports IORING\_\allowbreak{}OP\_\allowbreak{}WAITID (50, since 6.7): waitid without the syscall \\
\pid{io_uring.opcode.futex_wait} & seccomp & -- & 38/61 & IORING\_\allowbreak{}REGISTER\_\allowbreak{}PROBE reports IORING\_\allowbreak{}OP\_\allowbreak{}FUTEX\_\allowbreak{}WAIT (51, since 6.7): futex wait without futex(2) \\
\pid{io_uring.opcode.futex_wake} & seccomp & -- & 38/61 & IORING\_\allowbreak{}REGISTER\_\allowbreak{}PROBE reports IORING\_\allowbreak{}OP\_\allowbreak{}FUTEX\_\allowbreak{}WAKE (52, since 6.7): futex wake without futex(2) \\
\pid{io_uring.opcode.futex_waitv} & seccomp & -- & 38/61 & IORING\_\allowbreak{}REGISTER\_\allowbreak{}PROBE reports IORING\_\allowbreak{}OP\_\allowbreak{}FUTEX\_\allowbreak{}WAITV (53, since 6.7): futex\_\allowbreak{}waitv without the syscall \\
\pid{io_uring.opcode.fixed_fd_install} & seccomp & -- & 38/61 & IORING\_\allowbreak{}REGISTER\_\allowbreak{}PROBE reports IORING\_\allowbreak{}OP\_\allowbreak{}FIXED\_\allowbreak{}FD\_\allowbreak{}INSTALL (54, since 6.8): a ring-registered descriptor installed into the process's table \\
\pid{io_uring.opcode.ftruncate} & seccomp & -- & 38/61 & IORING\_\allowbreak{}REGISTER\_\allowbreak{}PROBE reports IORING\_\allowbreak{}OP\_\allowbreak{}FTRUNCATE (55, since 6.9): ftruncate without the syscall \\
\pid{io_uring.opcode.bind} & seccomp & -- & 38/61 & IORING\_\allowbreak{}REGISTER\_\allowbreak{}PROBE reports IORING\_\allowbreak{}OP\_\allowbreak{}BIND (56, since 6.11): bind without bind(2) \\
\pid{io_uring.opcode.listen} & seccomp & -- & 38/61 & IORING\_\allowbreak{}REGISTER\_\allowbreak{}PROBE reports IORING\_\allowbreak{}OP\_\allowbreak{}LISTEN (57, since 6.11): listen without listen(2) \\
\multicolumn{5}{@{}l}{\textbf{\texttt{mount}} (38)}\\
\pid{mount.reach} & seccomp & 2 & -- & mount(NULL, "/caliper-no-such-dir", "proc", 0, NULL): guaranteed ENOENT \\
\pid{mount.type.proc} & undec. & -- & 19 & /proc/filesystems lists proc: procfs, mountable anywhere the capability allows \\
\pid{mount.type.sysfs} & undec. & -- & 19 & /proc/filesystems lists sysfs: sysfs \\
\pid{mount.type.tmpfs} & undec. & -- & 19 & /proc/filesystems lists tmpfs: tmpfs \\
\pid{mount.type.devtmpfs} & undec. & -- & 19 & /proc/filesystems lists devtmpfs: the kernel's own /dev \\
\pid{mount.type.devpts} & undec. & -- & 19 & /proc/filesystems lists devpts: pseudoterminals \\
\pid{mount.type.overlay} & undec. & -- & 19 & /proc/filesystems lists overlay: overlayfs, every container's root \\
\pid{mount.type.fuse} & undec. & -- & 19 & /proc/filesystems lists fuse: user-space filesystems \\
\pid{mount.type.fusectl} & undec. & -- & 19 & /proc/filesystems lists fusectl: the FUSE control filesystem \\
\pid{mount.type.cgroup2} & undec. & -- & 19 & /proc/filesystems lists cgroup2: the unified cgroup hierarchy \\
\pid{mount.type.cgroup} & undec. & -- & 19 & /proc/filesystems lists cgroup: the v1 cgroup hierarchy \\
\pid{mount.type.bpf} & undec. & -- & 19 & /proc/filesystems lists bpf: pinned BPF objects \\
\pid{mount.type.debugfs} & undec. & -- & 19 & /proc/filesystems lists debugfs: kernel debugging interfaces \\
\pid{mount.type.tracefs} & undec. & -- & 19 & /proc/filesystems lists tracefs: ftrace \\
\pid{mount.type.configfs} & undec. & -- & 19 & /proc/filesystems lists configfs: kernel object configuration \\
\pid{mount.type.securityfs} & undec. & -- & 19 & /proc/filesystems lists securityfs: LSM interfaces \\
\pid{mount.type.efivarfs} & undec. & -- & 19 & /proc/filesystems lists efivarfs: UEFI variables \\
\pid{mount.type.pstore} & undec. & -- & 19 & /proc/filesystems lists pstore: persistent storage for crash logs \\
\pid{mount.type.binfmt_misc} & undec. & -- & 19 & /proc/filesystems lists binfmt\_\allowbreak{}misc: interpreter registration \\
\pid{mount.type.hugetlbfs} & undec. & -- & 19 & /proc/filesystems lists hugetlbfs: huge pages \\
\pid{mount.type.mqueue} & undec. & -- & 19 & /proc/filesystems lists mqueue: POSIX message queues \\
\pid{mount.type.autofs} & undec. & -- & 19 & /proc/filesystems lists autofs: automounter \\
\pid{mount.type.ecryptfs} & undec. & -- & 19 & /proc/filesystems lists ecryptfs: stacked encryption \\
\pid{mount.type.squashfs} & undec. & -- & 19 & /proc/filesystems lists squashfs: compressed read-only images \\
\pid{mount.type.iso9660} & undec. & -- & 19 & /proc/filesystems lists iso9660: optical media images \\
\pid{mount.type.vfat} & undec. & -- & 19 & /proc/filesystems lists vfat: FAT \\
\pid{mount.type.ext4} & undec. & -- & 19 & /proc/filesystems lists ext4: ext4 \\
\pid{mount.type.btrfs} & undec. & -- & 19 & /proc/filesystems lists btrfs: btrfs \\
\pid{mount.type.xfs} & undec. & -- & 19 & /proc/filesystems lists xfs: xfs \\
\pid{mount.type.nfs} & undec. & -- & 19 & /proc/filesystems lists nfs: NFS v3 \\
\pid{mount.type.nfs4} & undec. & -- & 19 & /proc/filesystems lists nfs4: NFS v4 \\
\pid{mount.type.cifs} & undec. & -- & 19 & /proc/filesystems lists cifs: SMB \\
\pid{mount.type.9p} & undec. & -- & 19 & /proc/filesystems lists 9p: Plan 9 transport (virtio-9p) \\
\pid{mount.type.virtiofs} & undec. & -- & 19 & /proc/filesystems lists virtiofs: virtio-fs \\
\pid{mount.type.ntfs3} & undec. & -- & 19 & /proc/filesystems lists ntfs3: NTFS, the in-kernel driver \\
\pid{mount.fs.proc} & undec. & -- & -- & mount(none, /tmp/caliper-\textless{}run\textgreater{}-mount, proc): a second procfs, unmounted by the child \\
\pid{mount.fs.overlay}$^{\dagger}$ & undec. & -- & -- & mount(none, /tmp/caliper-\textless{}run\textgreater{}-mount, overlay, lowerdir=a:b): a read-only overlay, unmounted by the child \\
\pid{mount.fs.fuse}$^{\dagger}$ & undec. & 22 & -- & mount(none, /tmp/caliper-\textless{}run\textgreater{}-mount, fuse): fuse without a descriptor, guaranteed EINVAL past the capability \\
\multicolumn{5}{@{}l}{\textbf{\texttt{clone}} (19)}\\
\pid{unshare.flags.invalid} & seccomp & 22 & -- & unshare(0x1): an invalid flag set, guaranteed EINVAL \\
\pid{unshare.flags.newuser} & sec./LSM & -- & -- & unshare(CLONE\_\allowbreak{}NEWUSER): user namespace creation without host privilege \\
\pid{unshare.flags.newns} & undec. & -- & -- & unshare(CLONE\_\allowbreak{}NEWNS): mount namespace creation, CAP\_\allowbreak{}SYS\_\allowbreak{}ADMIN \\
\pid{unshare.flags.newnet} & undec. & -- & -- & unshare(CLONE\_\allowbreak{}NEWNET): network namespace creation, CAP\_\allowbreak{}SYS\_\allowbreak{}ADMIN \\
\pid{unshare.flags.newpid} & undec. & -- & -- & unshare(CLONE\_\allowbreak{}NEWPID): PID namespace creation, CAP\_\allowbreak{}SYS\_\allowbreak{}ADMIN \\
\pid{unshare.flags.newipc} & undec. & -- & -- & unshare(CLONE\_\allowbreak{}NEWIPC): IPC namespace creation, CAP\_\allowbreak{}SYS\_\allowbreak{}ADMIN \\
\pid{unshare.flags.newuts} & undec. & -- & -- & unshare(CLONE\_\allowbreak{}NEWUTS): UTS namespace creation, CAP\_\allowbreak{}SYS\_\allowbreak{}ADMIN \\
\pid{unshare.flags.newcgroup} & undec. & -- & -- & unshare(CLONE\_\allowbreak{}NEWCGROUP): cgroup namespace creation, CAP\_\allowbreak{}SYS\_\allowbreak{}ADMIN \\
\pid{unshare.flags.newtime} & undec. & -- & -- & unshare(CLONE\_\allowbreak{}NEWTIME): time namespace creation, CAP\_\allowbreak{}SYS\_\allowbreak{}ADMIN \\
\pid{unshare.flags.newuser_newnet} & sec./LSM & -- & -- & unshare(CLONE\_\allowbreak{}NEWUSER \textbar{} CLONE\_\allowbreak{}NEWNET): user+net namespace creation escalation path \\
\pid{clone.flags.newuser} & sec./LSM & -- & -- & clone(CLONE\_\allowbreak{}NEWUSER \textbar{} SIGCHLD): process creation in a new user namespace \\
\pid{clone.flags.newns} & undec. & -- & -- & clone(CLONE\_\allowbreak{}NEWNS \textbar{} SIGCHLD): process creation in a new mount namespace, CAP\_\allowbreak{}SYS\_\allowbreak{}ADMIN \\
\pid{clone.flags.newnet} & undec. & -- & -- & clone(CLONE\_\allowbreak{}NEWNET \textbar{} SIGCHLD): process creation in a new net namespace, CAP\_\allowbreak{}SYS\_\allowbreak{}ADMIN \\
\pid{clone.flags.newpid} & undec. & -- & -- & clone(CLONE\_\allowbreak{}NEWPID \textbar{} SIGCHLD): process creation in a new PID namespace, CAP\_\allowbreak{}SYS\_\allowbreak{}ADMIN \\
\pid{clone.flags.newipc} & undec. & -- & -- & clone(CLONE\_\allowbreak{}NEWIPC \textbar{} SIGCHLD): process creation in a new IPC namespace, CAP\_\allowbreak{}SYS\_\allowbreak{}ADMIN \\
\pid{clone.flags.newuts} & undec. & -- & -- & clone(CLONE\_\allowbreak{}NEWUTS \textbar{} SIGCHLD): process creation in a new UTS namespace, CAP\_\allowbreak{}SYS\_\allowbreak{}ADMIN \\
\pid{clone.flags.newcgroup} & undec. & -- & -- & clone(CLONE\_\allowbreak{}NEWCGROUP \textbar{} SIGCHLD): process creation in a new cgroup namespace, CAP\_\allowbreak{}SYS\_\allowbreak{}ADMIN \\
\pid{clone.flags.newuser_newnet} & sec./LSM & -- & -- & clone(CLONE\_\allowbreak{}NEWUSER \textbar{} CLONE\_\allowbreak{}NEWNET \textbar{} SIGCHLD): user+net namespace process creation escalation path \\
\pid{clone3.args.short} & seccomp & 22 & 38 & clone3(NULL, 0): a size below CLONE\_\allowbreak{}ARGS\_\allowbreak{}SIZE\_\allowbreak{}VER0, guaranteed EINVAL \\
\multicolumn{5}{@{}l}{\textbf{\texttt{path}} (17)}\\
\pid{path.masked.proc_kcore} & sec./LSM & -- & 2 & openat(O\_\allowbreak{}PATH) + fstatfs: mask (ENODATA), read-only mount (EROFS), or reachable path \\
\pid{path.masked.proc_keys} & sec./LSM & -- & 2 & openat(O\_\allowbreak{}PATH) + fstatfs: mask (ENODATA), read-only mount (EROFS), or reachable path \\
\pid{path.masked.proc_latency_stats} & sec./LSM & -- & 2 & openat(O\_\allowbreak{}PATH) + fstatfs: mask (ENODATA), read-only mount (EROFS), or reachable path \\
\pid{path.masked.proc_timer_list} & sec./LSM & -- & 2 & openat(O\_\allowbreak{}PATH) + fstatfs: mask (ENODATA), read-only mount (EROFS), or reachable path \\
\pid{path.masked.proc_timer_stats} & sec./LSM & -- & 2 & openat(O\_\allowbreak{}PATH) + fstatfs: mask (ENODATA), read-only mount (EROFS), or reachable path \\
\pid{path.masked.proc_sched_debug} & sec./LSM & -- & 2 & openat(O\_\allowbreak{}PATH) + fstatfs: mask (ENODATA), read-only mount (EROFS), or reachable path \\
\pid{path.masked.proc_acpi} & sec./LSM & -- & 2 & openat(O\_\allowbreak{}PATH) + fstatfs: mask (ENODATA), read-only mount (EROFS), or reachable path \\
\pid{path.masked.proc_asound} & sec./LSM & -- & 2 & openat(O\_\allowbreak{}PATH) + fstatfs: mask (ENODATA), read-only mount (EROFS), or reachable path \\
\pid{path.masked.proc_scsi} & sec./LSM & -- & 2 & openat(O\_\allowbreak{}PATH) + fstatfs: mask (ENODATA), read-only mount (EROFS), or reachable path \\
\pid{path.masked.sys_firmware} & sec./LSM & -- & 2 & openat(O\_\allowbreak{}PATH) + fstatfs: mask (ENODATA), read-only mount (EROFS), or reachable path \\
\pid{path.masked.sys_powercap} & sec./LSM & -- & 2 & openat(O\_\allowbreak{}PATH) + fstatfs: mask (ENODATA), read-only mount (EROFS), or reachable path \\
\pid{path.readonly.proc_bus} & sec./LSM & -- & 2 & openat(O\_\allowbreak{}PATH) + fstatfs: mask (ENODATA), read-only mount (EROFS), or reachable path \\
\pid{path.readonly.proc_fs} & sec./LSM & -- & 2 & openat(O\_\allowbreak{}PATH) + fstatfs: mask (ENODATA), read-only mount (EROFS), or reachable path \\
\pid{path.readonly.proc_irq} & sec./LSM & -- & 2 & openat(O\_\allowbreak{}PATH) + fstatfs: mask (ENODATA), read-only mount (EROFS), or reachable path \\
\pid{path.readonly.proc_sys} & sec./LSM & -- & 2 & openat(O\_\allowbreak{}PATH) + fstatfs: mask (ENODATA), read-only mount (EROFS), or reachable path \\
\pid{path.readonly.proc_sysrq_trigger} & sec./LSM & -- & 2 & openat(O\_\allowbreak{}PATH) + fstatfs: mask (ENODATA), read-only mount (EROFS), or reachable path \\
\pid{path.control.proc_self_status} & sec./LSM & -- & 2 & openat(O\_\allowbreak{}PATH) + fstatfs: mask (ENODATA), read-only mount (EROFS), or reachable path \\
\multicolumn{5}{@{}l}{\textbf{\texttt{device}} (8)}\\
\pid{device.open.fuse}$^{\dagger}$ & undec. & -- & -- & /dev/fuse (10:229, typically 0666): read-only nonblocking open \\
\pid{device.open.net_tun}$^{\dagger}$ & undec. & -- & -- & /dev/net/tun (10:200, typically 0666): read-only nonblocking open \\
\pid{device.open.kvm}$^{\dagger}$ & undec. & -- & -- & /dev/kvm (10:232, typically 0660 root:kvm): DAC-gated read-only nonblocking open \\
\pid{device.open.kmsg} & undec. & -- & -- & /dev/kmsg (1:11, typically 0644): read-only nonblocking open \\
\pid{device.open.mem} & undec. & -- & -- & /dev/mem (1:1, typically 0640 root:kmem): DAC-gated read-only nonblocking open \\
\pid{device.open.port} & undec. & -- & -- & /dev/port (1:4, typically 0640 root:kmem): DAC-gated read-only nonblocking open \\
\pid{device.open.null} & undec. & -- & -- & /dev/null (1:3, typically 0666): control read-only nonblocking open \\
\pid{device.open.sda} & undec. & -- & -- & /dev/sda (8:0, typically 0660 root:disk): DAC-gated read-only nonblocking open \\
\multicolumn{5}{@{}l}{\textbf{\texttt{capability}} (8)}\\
\pid{capability.sys_time} & undec. & -- & -- & settimeofday(NULL, NULL): capability-only time check \\
\pid{capability.sys_boot} & undec. & 22 & -- & reboot(0, 0, LINUX\_\allowbreak{}REBOOT\_\allowbreak{}CMD\_\allowbreak{}CAD\_\allowbreak{}OFF, NULL): capability check \\
\pid{capability.sys_module} & undec. & 9 & -- & finit\_\allowbreak{}module(-1, "", 0): module capability check \\
\pid{capability.mknod} & undec. & -- & -- & mknodat(AT\_\allowbreak{}FDCWD, /tmp/caliper-mknod, S\_\allowbreak{}IFCHR\textbar{}0600, makedev(1,3)) \\
\pid{capability.sys_ptrace} & undec. & 14 & -- & process\_\allowbreak{}vm\_\allowbreak{}readv(getppid(), valid local iovec, invalid remote iovec): access check \\
\pid{capability.net_bind_service} & undec. & -- & -- & socket(AF\_\allowbreak{}INET) + bind(0.0.0.0:1): privileged-port check \\
\pid{capability.sys_nice} & undec. & -- & -- & setpriority(PRIO\_\allowbreak{}PROCESS, 0, -1): priority capability check \\
\pid{capability.sys_rawio} & undec. & -- & -- & iopl(3): x86 I/O privilege capability check (x86\_64 only) \\
\midrule\multicolumn{5}{@{}l}{\textbf{Deferred set: specified, never run}}\\
\multicolumn{5}{@{}l}{\textbf{\texttt{ioctl}} (3)}\\
\pid{ioctl.request.fioclex} & seccomp & 9 & -- & ioctl(-1, FIOCLEX, 0): the request every workload issues \\
\pid{ioctl.request.tiocsti} & seccomp & 9 & -- & ioctl(-1, TIOCSTI, NULL): terminal input injection \\
\pid{ioctl.request.tioclinux} & seccomp & 9 & -- & ioctl(-1, TIOCLINUX, NULL): virtual console control \\
\multicolumn{5}{@{}l}{\textbf{\texttt{bpf}} (5)}\\
\pid{bpf.cmd.map_create} & undec. & 7 & 38 & bpf(BPF\_\allowbreak{}MAP\_\allowbreak{}CREATE, attr, sizeof(union bpf\_\allowbreak{}attr) + 8) with a non-zero byte past the union \\
\pid{bpf.cmd.prog_load} & undec. & 7 & 38 & bpf(BPF\_\allowbreak{}PROG\_\allowbreak{}LOAD, attr, sizeof(union bpf\_\allowbreak{}attr) + 8) with a non-zero byte past the union \\
\pid{bpf.cmd.btf_load} & seccomp & 7 & 38 & bpf(BPF\_\allowbreak{}BTF\_\allowbreak{}LOAD, attr, sizeof(union bpf\_\allowbreak{}attr) + 8) with a non-zero byte past the union \\
\pid{bpf.prog_type.socket_filter} & undec. & 14 & 38 & bpf(BPF\_\allowbreak{}PROG\_\allowbreak{}LOAD, \{prog\_\allowbreak{}type: SOCKET\_\allowbreak{}FILTER, insn\_\allowbreak{}cnt: 1, insns: NULL\}, sizeof(union bpf\_\allowbreak{}attr)): the type unprivileged loaders may use \\
\pid{bpf.prog_type.kprobe} & undec. & 14 & 38 & bpf(BPF\_\allowbreak{}PROG\_\allowbreak{}LOAD, \{prog\_\allowbreak{}type: KPROBE, insn\_\allowbreak{}cnt: 1, insns: NULL\}, sizeof(union bpf\_\allowbreak{}attr)): a tracing type \\
\multicolumn{5}{@{}l}{\textbf{\texttt{keyctl}} (4)}\\
\pid{keyctl.op.read} & seccomp & 126 & 38 & keyctl(KEYCTL\_\allowbreak{}READ, 0x7fffffff, NULL, 0) \\
\pid{keyctl.op.out_of_range} & seccomp & 95 & 38 & keyctl(0x7fffffff, 0, 0, 0, 0): an operation no kernel defines \\
\pid{keyctl.add_key} & seccomp & 126 & 38 & add\_\allowbreak{}key("user", "caliper", NULL, 0, 0x7fffffff) \\
\pid{keyctl.request_key} & seccomp & 126 & 38 & request\_\allowbreak{}key("user", "caliper", NULL, 0x7fffffff) \\
\multicolumn{5}{@{}l}{\textbf{\texttt{prctl}} (3)}\\
\pid{prctl.option.out_of_range} & sec./LSM & 22 & -- & prctl(0x7fffffff, 0, 0, 0, 0): an option no kernel defines \\
\pid{prctl.option.set_mm} & sec./LSM & 22 & -- & prctl(PR\_\allowbreak{}SET\_\allowbreak{}MM, PR\_\allowbreak{}SET\_\allowbreak{}MM\_\allowbreak{}START\_\allowbreak{}CODE, 0, 0, 1): a non-zero arg5 \\
\pid{prctl.option.set_dumpable} & sec./LSM & 22 & -- & prctl(PR\_\allowbreak{}SET\_\allowbreak{}DUMPABLE, 2, 0, 0, 0): SUID\_\allowbreak{}DUMP\_\allowbreak{}ROOT, which prctl refuses \\
\multicolumn{5}{@{}l}{\textbf{\texttt{perf\_event\_open}} (3)}\\
\pid{perf_event_open.flags.invalid} & seccomp & 22 & 38 & perf\_\allowbreak{}event\_\allowbreak{}open(\&attr, 0, -1, -1, \textasciitilde{}0UL) \\
\pid{perf_event_open.attr.bad_size} & seccomp & 7 & 38 & perf\_\allowbreak{}event\_\allowbreak{}open(\{size: 1\}, 0, -1, -1, 0) \\
\pid{perf_event_open.cpu_wide} & sec./LSM & -- & 38 & perf\_\allowbreak{}event\_\allowbreak{}open(\{type: SOFTWARE, config: CPU\_\allowbreak{}CLOCK\}, -1, 0, -1, PERF\_\allowbreak{}FLAG\_\allowbreak{}FD\_\allowbreak{}CLOEXEC): every task on CPU 0 \\
\multicolumn{5}{@{}l}{\textbf{\texttt{userfaultfd}} (3)}\\
\pid{userfaultfd.flags.invalid} & seccomp & 22 & 38 & userfaultfd(UFFD\_\allowbreak{}USER\_\allowbreak{}MODE\_\allowbreak{}ONLY \textbar{} 2) \\
\pid{userfaultfd.user_mode_only} & sec./LSM & -- & 38 & userfaultfd(UFFD\_\allowbreak{}USER\_\allowbreak{}MODE\_\allowbreak{}ONLY \textbar{} O\_\allowbreak{}CLOEXEC) \\
\pid{userfaultfd.kernel_faults} & undec. & -- & 38 & userfaultfd(O\_\allowbreak{}CLOEXEC): kernel-mode faults \\
\multicolumn{5}{@{}l}{\textbf{\texttt{open\_by\_handle\_at}} (2)}\\
\pid{open_by_handle_at.open} & undec. & 14 & 38 & open\_\allowbreak{}by\_\allowbreak{}handle\_\allowbreak{}at(AT\_\allowbreak{}FDCWD, NULL, O\_\allowbreak{}RDONLY) \\
\pid{open_by_handle_at.name_to_handle_at} & seccomp & 22 & 38 & name\_\allowbreak{}to\_\allowbreak{}handle\_\allowbreak{}at(AT\_\allowbreak{}FDCWD, "/", NULL, NULL, \textasciitilde{}0) \\
\multicolumn{5}{@{}l}{\textbf{\texttt{module}} (2)}\\
\pid{module.init_module} & undec. & 8 & 38 & init\_\allowbreak{}module(NULL, 0, "") \\
\pid{module.delete_module} & undec. & 2 & 38 & delete\_\allowbreak{}module("caliper\_\allowbreak{}absent", O\_\allowbreak{}NONBLOCK) \\
\multicolumn{5}{@{}l}{\textbf{\texttt{seccomp\_notify}} (3)}\\
\pid{seccomp_notify.new_listener} & seccomp & 14 & 38 & seccomp(SECCOMP\_\allowbreak{}SET\_\allowbreak{}MODE\_\allowbreak{}FILTER, SECCOMP\_\allowbreak{}FILTER\_\allowbreak{}FLAG\_\allowbreak{}NEW\_\allowbreak{}LISTENER, NULL) \\
\pid{seccomp_notify.get_notif_sizes} & seccomp & 14 & 38 & seccomp(SECCOMP\_\allowbreak{}GET\_\allowbreak{}NOTIF\_\allowbreak{}SIZES, 0, NULL) \\
\pid{seccomp_notify.ioctl_recv} & seccomp & 9 & 38 & ioctl(-1, SECCOMP\_\allowbreak{}IOCTL\_\allowbreak{}NOTIF\_\allowbreak{}RECV, NULL) \\
\end{longtable}\endgroup

